\documentclass[11pt]{article}
\usepackage[T1]{fontenc}
\usepackage{lmodern}
\usepackage{amsmath,amssymb,amsthm}
\usepackage[margin=1in]{geometry}
\usepackage{microtype}
\usepackage[colorlinks=true,linkcolor=blue,citecolor=blue,urlcolor=blue]{hyperref}

\newtheorem{theorem}{Theorem}
\newtheorem{lemma}[theorem]{Lemma}
\newcommand{\C}{\mathbb C}
\newcommand{\R}{\mathbb R}
\newcommand{\Z}{\mathbb Z}

\title{Relative density under fixed-parameter extrapolation}
\author{Kaiwen Jin\\\small\texttt{yui.ui.siki@gmail.com}}
\date{October 14, 2026}

\begin{document}
\maketitle

\begin{abstract}
Starting from $0$ and $1$, repeatedly apply the operation
$(a,b)\mapsto(1-\lambda)a+\lambda b$, with a fixed complex parameter
$\lambda$. We prove that the generated set has bounded holes in the real
line for every real $\lambda\notin[0,1]$, and in the complex plane for every
nonreal $\lambda$. This answers the relative-density question of Fenner,
Green and Homer. Their characterization by integer polynomials allows a
single affine image of a relatively dense set of finite digit expansions
to be placed inside the generated set.
\end{abstract}

\section{The relative-density question}

For $\lambda\in\C$, let $Q_\lambda$ be the smallest subset of $\C$
containing $0,1$ and closed under
\[
 a\star_\lambda b=(1-\lambda)a+\lambda b.
\]
Only finitely many applications of this operation are allowed; no
topological closure is taken. Such sets occur in the study of aperiodic
point patterns. Relative density asks whether their holes have bounded size.
Fenner, Green and Homer established this property for several algebraic
parameters and asked whether it holds throughout the following two domains
\cite[Section~6 and Open Question~11.3]{FGH}.

For $\lambda\in\C\setminus[0,1]$, set
\[
 X_\lambda=
 \begin{cases}
  \R,&\lambda\in\R\setminus[0,1],\\
  \C,&\lambda\in\C\setminus\R.
 \end{cases}
\]

\begin{theorem}\label{thm:main}
For every $\lambda\in\C\setminus[0,1]$, there is a finite constant
$R_\lambda$ such that
\[
 \forall z\in X_\lambda\quad\exists q\in Q_\lambda:
 \qquad |z-q|\le R_\lambda.
\]
\end{theorem}

The radius may depend on the parameter. The proof uses integer polynomials
that are small on $[0,1]$ and large at $\lambda$. Finite sums of their powers
can approximate every point with a fixed error, while remaining uniformly
bounded on the interval. We then use the bounded-polynomial embedding
observation in \cite[Lemma~6.19 and the following remark]{FGH} to obtain
points of $Q_\lambda$ at one fixed scale.

\section{Proof}

We use the following consequence of the polynomial characterization of
Fenner, Green and Homer \cite[Theorem~3.6, Fact~3.3 and Corollary~3.8]{FGH}:
\begin{equation}\label{eq:characterization}
 p\in\Z[x],\qquad 0<p(t)<1\ \text{for every }0<t<1
 \quad\Longrightarrow\quad p(\lambda)\in Q_\lambda.
\end{equation}
The conclusion concerns the generated set itself.

\begin{lemma}\label{lem:separation}
For every $\lambda\in\C\setminus[0,1]$, there is $h\in\Z[x]$ such that
\[
 \sup_{0\le t\le1}|h(t)|\le\frac14,
 \qquad |h(\lambda)|>1.
\]
\end{lemma}

\begin{proof}
Let $C_0(u)=1$, $C_1(u)=u$, and
$C_{n+1}(u)=2uC_n(u)-C_{n-1}(u)$. These are integer polynomials. The
recurrence gives, for every $r\ne0$,
\begin{equation}\label{eq:chebyshev}
 C_n\!\left(\frac{r+r^{-1}}2\right)=\frac{r^n+r^{-n}}2.
\end{equation}
For $u\in[-1,1]$, choose $r=u+i\sqrt{1-u^2}$, which has modulus one.
Equation~\eqref{eq:chebyshev} gives $|C_n(u)|\le1$.

Put $w=2\lambda-1\notin[-1,1]$. The roots of $r^2-2wr+1=0$ are
nonzero and reciprocal. If both had modulus one, then
$w=(r+\overline r)/2\in[-1,1]$. Thus one root has $|r|>1$, and
\[
 |C_n(w)|\ge\frac{|r|^n-|r|^{-n}}2\longrightarrow\infty.
\]
Writing $s(x)=x(1-x)$, choose $n$ so that
$|s(\lambda)C_n(2\lambda-1)|>1$, which is possible since
$s(\lambda)\ne0$. Then $h(x)=s(x)C_n(2x-1)$ has the required properties.
\end{proof}

\begin{lemma}\label{lem:digits}
Fix $\lambda\in\C\setminus[0,1]$, and put
\[
 L=\begin{cases}\Z,&X_\lambda=\R,\\
                  \Z+\lambda\Z,&X_\lambda=\C.
    \end{cases}
\]
If $\beta\in X_\lambda$ and $|\beta|>1$, there are a finite set
$D\subset L$ containing $0$ and a finite $\rho>0$ such that
\[
 S=\left\{\sum_{j=0}^m d_j\beta^j:
                 m\ge0,\ d_j\in D\right\}
\]
is $\rho$-relatively dense in $X_\lambda$.
\end{lemma}

\begin{proof}
Take $\delta=1/2$ in the real case and
$\delta=(1+|\lambda|)/2$ in the nonreal case. Every point of $X_\lambda$
is within distance $\delta$ of $L$. In the second case, write the point as
$a+b\lambda$ with $a,b\in\R$ and round both coefficients to integers.
This is possible because $\operatorname{Im}\lambda\ne0$.
Also, bounded subsets of $L$ are finite: if $|m+n\lambda|\le A$, then
\[
 |n|\le\frac A{|\operatorname{Im}\lambda|},\qquad
 |m|\le A+\frac{A|\operatorname{Re}\lambda|}{|\operatorname{Im}\lambda|}.
\]
In the real case this finiteness follows from $|m|\le A$.

Set $\rho=|\beta|\delta$,
$K=\{y\in X_\lambda:|y|\le\rho\}$ and
$D=L\cap\{d:|d|\le\rho+\delta\}$. For each $y\in K$, choose $d\in L$
with $|y-d|\le\delta$. Then $d\in D$ and $\beta(y-d)\in K$. Thus
\begin{equation}\label{eq:cover}
 K\subseteq D+\beta^{-1}K.
\end{equation}

Given $z\in X_\lambda$, choose $N\ge0$ with
$y_0=\beta^{-N}z\in K$. If $N=0$, the element $0\in S$ is sufficient.
Otherwise apply \eqref{eq:cover} finitely many times to choose $d_j\in D$
and $y_{j+1}=\beta(y_j-d_j)\in K$ for $j=0,\ldots,N-1$. This gives
\begin{equation}\label{eq:finite-expansion}
 z=\sum_{j=0}^{N-1}\beta^{N-j}d_j+y_N.
\end{equation}
The sum is in $S$, with zero as its constant digit, and $|y_N|\le\rho$.
\end{proof}

\begin{proof}[Proof of Theorem~\ref{thm:main}]
Fix $\lambda\in\C\setminus[0,1]$. Choose $h$ from
Lemma~\ref{lem:separation} and put $\beta=h(\lambda)$. When $\lambda$ is
real, so is $\beta$, and hence Lemma~\ref{lem:digits} applies in both
domains. Let $D,S,\rho$ be as in that lemma.

For each $d\in D$, choose an integer polynomial $e_d$ satisfying
$e_d(\lambda)=d$: use the constant polynomial $d$ in the real case,
and $e_d(x)=a+bx$ when $d=a+b\lambda$, $a,b\in\Z$, in the nonreal case.
Define
\[
 M=\max\!\left(1,\max_{d\in D}\max\{|e_d(0)|,|e_d(1)|\}\right),
 \qquad B=\frac{4M}{3}.
\]
For $0\le t\le1$, the affine identity
$e_d(t)=(1-t)e_d(0)+te_d(1)$ gives $|e_d(t)|\le M$. Consequently every
finite digit polynomial
\[
 f(x)=\sum_{j=0}^m e_{d_j}(x)h(x)^j
\]
has integer coefficients and satisfies the bound
\begin{equation}\label{eq:uniform-bound}
 |f(t)|\le M\sum_{j=0}^m4^{-j}\le B
 \qquad(0\le t\le1).
\end{equation}
Their evaluations at $\lambda$ exhaust $S$.

Choose an integer $\ell>B$ and an integer $k\ge1$ with
$(\ell+B)4^{-k}<1$. These choices are the same for every digit polynomial.
With $s(x)=x(1-x)$, set
\[
 p_f(x)=s(x)^k(\ell+f(x))\in\Z[x].
\]
For every $0<t<1$, \eqref{eq:uniform-bound} implies
\[
 0<s(t)^k(\ell-B)\le p_f(t)\le4^{-k}(\ell+B)<1.
\]
The characterization \eqref{eq:characterization} therefore yields
\begin{equation}\label{eq:embedding}
 c(\ell+S)\subseteq Q_\lambda,
 \qquad c=s(\lambda)^k\ne0.
\end{equation}

Finally, for any $z\in X_\lambda$, choose $v\in S$ with
$|z/c-\ell-v|\le\rho$, and put $q=c(\ell+v)\in Q_\lambda$. Then
\[
 |z-q|\le |c|\rho.
\]
Thus $R_\lambda=|\lambda(1-\lambda)|^k\rho$ is a finite covering radius.
The quantities $h,D,\rho,\ell,k$ are fixed before $z$ is chosen, and every
approximant in \eqref{eq:embedding} comes from a finite digit polynomial.
\end{proof}

\paragraph{AI-use statement.}
OpenAI Codex was used to develop the mathematical argument, check the cited
source, and prepare this manuscript.

\begin{thebibliography}{1}
\bibitem{FGH}
S.~Fenner, F.~Green and S.~Homer,
\emph{Fixed-Parameter Extrapolation and Aperiodic Order},
Discrete \& Computational Geometry \textbf{76} (2026), 1--74.
\href{https://doi.org/10.1007/s00454-025-00816-4}{doi:10.1007/s00454-025-00816-4}.
\end{thebibliography}
\end{document}
